% AI-assisted manuscript source; responsible author: Carptopus.
\documentclass[11pt]{article}
\usepackage[a4paper,margin=29mm]{geometry}
\usepackage[T1]{fontenc}
\usepackage{lmodern}
\usepackage{microtype}
\usepackage{amsmath,amssymb,amsthm,mathtools}
\usepackage{enumitem}
\usepackage{fvextra}
\usepackage{xcolor}
\usepackage[colorlinks=true,linkcolor=blue!55!black,citecolor=blue!55!black,urlcolor=blue!65!black]{hyperref}
\usepackage{fancyhdr}

\newtheorem{theorem}{Theorem}[section]
\newtheorem{lemma}[theorem]{Lemma}
\newtheorem{proposition}[theorem]{Proposition}
\newtheorem{corollary}[theorem]{Corollary}
\theoremstyle{remark}
\newtheorem{remark}[theorem]{Remark}

\pagestyle{fancy}
\fancyhf{}
\fancyhead[L]{Carptopus}
\fancyhead[R]{The nSSP for root-loop trees}
\fancyfoot[C]{\thepage}
\setlength{\headheight}{14pt}
\setlength{\parindent}{0pt}
\setlength{\parskip}{0.55em}
\setlength{\emergencystretch}{2.5em}
\setlist{nosep,leftmargin=*}

\title{A recursive non-symmetric strong spectral criterion\\for root-loop tree matrices}
\author{Carptopus\\\href{mailto:carptopus@163.com}{carptopus@163.com}}
\date{22 August 2026}

\hypersetup{
  pdftitle={A recursive non-symmetric strong spectral criterion for root-loop tree matrices},
  pdfauthor={Carptopus},
  pdfsubject={A recursive coprimality criterion for the nSSP of root-loop bidirected tree matrices},
  pdfkeywords={nSSP, rooted tree matrix, inverse eigenvalue problem, spectral criterion, coprime polynomials}
}

\begin{document}
\maketitle

\begin{center}
\small\itshape
Public Beta v0.2. The theorem chain has passed a writing-admission
proof audit, exact computational calibration, and a manuscript-level zero-trust audit.
External mathematical review and journal peer review are pending. Priority statements
are qualified by the public-literature boundary stated in the manuscript.
\end{center}
\vspace{0.8em}

\begin{abstract}

Let $(T,r)$ be a finite rooted tree. We consider arbitrary real matrices whose nonzero off-diagonal positions are exactly the two directed entries on every edge of $T$, whose root entry is a single nonzero loop, and whose remaining diagonal entries are zero. No symmetry, sign, real-spectrum, diagonalizability, or square-free-spectrum assumption is imposed. For every vertex $v$ and child $c$ of $v$, let $F_c(x)$ be the characteristic polynomial of the complete rooted subtree below $c$.

We prove that the full matrix has the non-symmetric strong spectral property if and only if, at every vertex, the child-subtree polynomials are pairwise coprime. The sufficient direction combines a recursive PBH cyclicity argument with an exact infinitesimal recovery of all directed edge products from the characteristic polynomial. The converse is constructive: a deepest collision between two child-subtree spectra produces a nonzero real pattern-zero centralizer, including for a zero, nonreal, or repeated common root. The theorem strictly extends the root-loop spider criterion to arbitrarily many nested branch vertices. It also gives a topological obstruction: a positive instance must satisfy the classical nearly-even-branching recursion. Exact rational and modular computations are included only as reproducible falsification controls and do not replace the proof.

\end{abstract}

\section{Introduction}


The non-symmetric strong spectral property, abbreviated nSSP, is a transversality condition for a real matrix and its zero--nonzero or sign pattern \cite{ref1,ref2,ref11}. It is useful in bifurcation and superpattern arguments for non-symmetric inverse eigenvalue problems. A pattern-level question asks whether some or every matrix in a qualitative class has the property. The present question has different quantifiers: it asks for an exact criterion for each fixed matrix in one sparse infinite family.

The family consists of bidirected tree matrices with exactly one nonzero diagonal position, at a chosen root. The two-arm fixed-matrix mechanism first appeared in the author's looped-double-path work \cite{ref12}, while deleting the root of a spider produces a direct sum of paths and leads to the pairwise-coprimality criterion in \cite{ref7}. For a general rooted tree, however, branch vertices occur recursively inside the deleted components. A root-level test is therefore insufficient: two grandchildren may share a spectral factor even when the root has only one child.

The main theorem shows that the spider mechanism persists, but only after being imposed at every vertex. The correct invariant is a hierarchy of characteristic polynomials of complete child subtrees. Pairwise coprimality at every sibling set is both necessary and sufficient.

Three points distinguish the statement from the classical Hermitian setting.

\begin{enumerate}
\item Directed edge products may have mixed signs, so the matrix need not be diagonally similar over $\mathbb R$ to a symmetric matrix.
\item A single child-subtree polynomial may have repeated or nonreal roots; only a factor shared by siblings is forbidden.
\item Failure is witnessed explicitly at the deepest spectral collision, rather than inferred only from a Jacobian rank loss.
\end{enumerate}


The recursive tree skeleton has important precedents. Duarte's construction and the 2013 Jacobian formulation use recursive strict interlacing for symmetric tree matrices \cite{ref3,ref8}. A skew-symmetric analogue connects the same recursion with nearly even branching (NEB) \cite{ref4}. These results cover the positive-product or skew-symmetric special mechanisms but do not state the mixed-product fixed-matrix nSSP equivalence proved below. The general characteristic-Jacobian-to-nSSP equivalence is also prior work \cite[Theorem 5.1]{ref1} and is used as a bridge, not claimed as a contribution.

\subsection*{Main theorem}


Let $(T,r)$ be a finite rooted tree and let $A$ be a real matrix such that:

\begin{itemize}
\item $A_{rr}=a\ne0$ and $A_{vv}=0$ for $v\ne r$;
\item for every tree edge $\{u,v\}$, both $A_{uv}$ and $A_{vu}$ are nonzero;
\item all remaining entries are zero.
\end{itemize}


For a vertex $v$, write $C(v)$ for its children. For every child $c\in C(v)$, let $B_c$ be the principal matrix on the complete rooted subtree below $c$, and set

\[
F_c(x)=\det(xI-B_c).
\]

Then

\[
A\text{ has the nSSP}
\quad\Longleftrightarrow\quad
\{F_c:c\in C(v)\}\text{ is pairwise coprime for every }v.
\tag{1.1}
\]

Equivalently, $A$ fails the nSSP exactly when two sibling subtrees share an eigenvalue over $\mathbb C$.

\subsection*{Contribution boundary}


The nSSP and bifurcation framework \cite{ref1,ref2}, its similarity-transversality formulation \cite{ref11}, the Duarte property \cite{ref3,ref8}, the skew-Duarte/NEB mechanism \cite{ref4}, PBH cyclicity, and the characteristic polynomial of a tree in terms of two-way edge products are prior tools. Matrices on the same rooted tree with the same root loop and the same edge products are related by real diagonal similarity; changing the directional splitting is therefore a gauge freedom, not an independent result.

The retained contribution is the exact synthesis into an arbitrary mixed-product fixed-tree criterion, including nonnormal and defective cases, together with a recursive recovery proof and an explicit real witness at an arbitrary-depth collision. This draft makes no absolute global-priority claim.

\section{Definitions and rooted-tree notation}


\subsection*{The nSSP}


For real matrices $A$ and $X$ of order $n$, let $A\circ X$ denote their Hadamard product. The matrix $A$ has the non-symmetric strong spectral property if

\[
A\circ X=0,
\qquad
AX^{\mathsf T}=X^{\mathsf T}A
\tag{2.1}
\]

imply $X=0$.

Equivalently, the tangent space to the zero--nonzero pattern class and the tangent space to the similarity orbit of $A$ span $M_n(\mathbb R)$ \cite{ref1,ref2}. We use (2.1) for the failure direction and the characteristic-coefficient Jacobian for the sufficient direction.

\subsection*{Rooted subtrees}


Orient the tree only for notation, away from $r$. For $v\ne r$, write $\pi(v)$ for its parent. Let $T_v$ be the complete descendant subtree rooted at $v$.

For every $v$, define $B_v$ to be the zero-diagonal matrix induced on $T_v$. When $v=r$, this means that the loop $a$ is temporarily suppressed. Put

\[
F_v(x)=\det(xI-B_v),
\qquad
Q_v(x)=\prod_{c\in C(v)}F_c(x).
\tag{2.2}
\]

The empty product is $1$. Thus a leaf has $F_v=x$ and $Q_v=1$.

For an edge from $v$ to its child $c$, write

\[
p_{vc}=A_{vc},
\qquad
q_{vc}=A_{cv},
\qquad
\beta_{vc}=p_{vc}q_{vc}\ne0.
\tag{2.3}
\]

The theorem's condition is

\[
\gcd(F_c,F_d)=1
\quad\text{for all }v\text{ and distinct }c,d\in C(v).
\tag{H}
\]

For a leaf or a vertex with one child, (H) is vacuous at that vertex.

\section{Polynomial recursion and edge-product reduction}


\begin{lemma}

For every vertex $v$,

\[
F_v
=xQ_v-
\sum_{c\in C(v)}
\beta_{vc}Q_c\prod_{\substack{d\in C(v)\\d\ne c}}F_d.
\tag{3.1}
\]

The characteristic polynomial of the full matrix is

\[
\chi_A(x)=F_r(x)-aQ_r(x).
\tag{3.2}
\]

\end{lemma}

\begin{proof}

A nonzero permutation term in the determinant of a zero-diagonal tree matrix is a product of disjoint transpositions along tree edges. Hence every directed edge enters only through $A_{uv}A_{vu}$. Expanding at $v$, the unmatched term contributes $xQ_v$. Matching $v$ to a child $c$ contributes $-\beta_{vc}$, the deleted-root polynomial $Q_c$ inside that child subtree, and the full polynomials of all sibling subtrees. This is (3.1). Restoring the root loop is a rank-one diagonal update whose cofactor is $Q_r$, giving (3.2).

\end{proof}

\begin{lemma}

The polynomial $F_v$ contains only powers of the same parity as $|T_v|$, while $Q_v$ contains only powers of the same parity as $|T_v|-1$.

\end{lemma}

\begin{proof}

Every determinant term of $B_v$ chooses a matching with $k$ edges and contributes a power $x^{|T_v|-2k}$. The statement for $Q_v$ follows either from the forest interpretation or from (2.2).

\end{proof}

\begin{remark}

The characteristic data depend only on $a$ and the edge products $\beta_e$. If two matrices on the same rooted tree have the same $a$ and the same $\beta_e$, choose a nonzero diagonal scaling recursively along the tree to match one directed entry on each edge; equality of the products then matches the reverse entry. Thus the matrices are real diagonally similar. The nSSP is invariant under this similarity because a pattern-zero centralizer is transported by the same diagonal conjugation.

\end{remark}

\section{Recursive cyclicity}


We work over $\mathbb C$ when discussing eigenvectors.

\begin{lemma}

Assume (H) holds at every vertex of $T_v$. Then $e_v$ is cyclic for both $B_v$ and $B_v^{\mathsf T}$.

\end{lemma}

\begin{proof}

We induct on $|T_v|$. The assertion is immediate for a leaf. Suppose a left eigenvector $y^{\mathsf T}B_v=\lambda y^{\mathsf T}$ has $y_v=0$. Its restriction to every child subtree is either zero or a left $\lambda$-eigenvector of the corresponding $B_c$. Pairwise coprimality allows $\lambda$ to occur in at most one child polynomial, so at most one child block is nonzero.

If the block below $c$ is nonzero, the induction hypothesis implies that its coordinate at $c$ is nonzero. The eigenvector equation at the parent $v$ then consists of one nonzero edge weight times this nonzero coordinate, a contradiction. Hence no left eigenvector is orthogonal to $e_v$. The PBH criterion \cite{ref10} gives cyclicity of $e_v$ for $B_v$.

Applying the same argument to $B_v^{\mathsf T}$ proves the second assertion, since transposition does not change any $F_c$.

\end{proof}

\begin{corollary}

Under (H), the root vector $e_r$ is cyclic for the full matrix $A$. In particular, the minimal polynomial of $A$ has degree $n$.

\end{corollary}

\begin{proof}

The argument of Lemma 4.1 is unchanged by the root loop when the root coordinate of the putative left eigenvector is zero.

The proof does not assume real or simple eigenvalues. Cyclicity permits a repeated characteristic root provided it belongs to one Jordan chain \cite{ref9}.

\end{proof}

\section{Infinitesimal recovery of the edge products}


Only the variables $\beta_e$ inside $T_v$ vary in this section. A dot denotes their first-order variation. We use the convention $\deg0=-\infty$.

\begin{lemma}

Assume (H) holds throughout $T_v$. If

\[
\dot F_v=0,
\qquad
\dot Q_v=0,
\tag{5.1}
\]

then every $\dot\beta_e$ in $T_v$ is zero.

\end{lemma}

\begin{proof}

We induct on the subtree order. A leaf has no edge variable. Differentiating $Q_v=\prod_cF_c$ gives

\[
0=\dot Q_v=
\sum_c\dot F_c\prod_{d\ne c}F_d.
\tag{5.2}
\]

Modulo $F_c$, the product of the siblings is invertible because of (H), so $F_c\mid\dot F_c$. The polynomial $F_c$ is monic and its coefficient of degree $\deg F_c-1$ is the fixed value $-\operatorname{tr}B_c=0$. Therefore

\[
\deg\dot F_c\le\deg F_c-2,
\]

and $\dot F_c=0$.

Differentiate (3.1) and use $\dot F_v=\dot Q_v=\dot F_c=0$:

\[
\sum_c
\left(\dot\beta_{vc}Q_c+\beta_{vc}\dot Q_c\right)
\prod_{d\ne c}F_d=0.
\tag{5.3}
\]

Reducing modulo $F_c$ yields

\[
F_c\mid H_c,
\qquad
H_c:=\dot\beta_{vc}Q_c+\beta_{vc}\dot Q_c.
\tag{5.4}
\]

Now $Q_c$ is monic of degree $\deg F_c-1$, and its next coefficient is again fixed at zero. Hence

\[
\deg\dot Q_c\le\deg F_c-3.
\]

Thus $\deg H_c<\deg F_c$, so $H_c=0$. Comparing its highest-degree coefficient gives $\dot\beta_{vc}=0$; since $\beta_{vc}\ne0$, it follows that $\dot Q_c=0$. Together with $\dot F_c=0$, the induction hypothesis recovers every descendant product.

For a leaf child, $Q_c=1$ and $\dot Q_c=0$; for one child, the same reductions hold without a sibling condition; empty products are $1$. These cases close all degree-bound endpoints.

\end{proof}

\section{Sufficiency: the characteristic Jacobian}


Let

\[
\Phi:(a,(\beta_e)_{e\in E(T)})
\longmapsto
\text{the }n\text{ nonleading coefficients of }\chi_A.
\tag{6.1}
\]

\begin{proposition}

Under (H), the Jacobian of $\Phi$ is nonsingular.

\end{proposition}

\begin{proof}

Suppose $\dot\chi_A=0$. By (3.2),

\[
0=\dot F_r-\dot a\,Q_r-a\dot Q_r.
\tag{6.2}
\]

Lemma 3.2 separates this identity into opposite parity parts:

\[
\dot F_r=0,
\qquad
\dot a\,Q_r+a\dot Q_r=0.
\tag{6.3}
\]

The polynomial $Q_r$ is monic while $\dot Q_r$ has degree at least two lower. Comparing leading coefficients gives $\dot a=0$. Since $a\ne0$, also $\dot Q_r=0$. Lemma 5.1 then gives $\dot\beta_e=0$ for every edge. The differential is injective between two $n$-dimensional spaces and is therefore nonsingular.

\end{proof}

\begin{proposition}

If (H) holds, then $A$ has the nSSP.

\end{proposition}

\begin{proof}

For one edge with directed entries $p,q\ne0$,

\[
d\beta=q\,dp+p\,dq,
\]

so the map from the original pattern variables to $(a,(\beta_e))$ has full row rank. Proposition 6.1 therefore implies that the full characteristic-coefficient Jacobian with respect to the nonzero entries of $A$ has row rank $n$.

Corollary 4.2 says that the minimal polynomial of $A$ has degree $n$. The general Jacobian--nSSP equivalence \cite[Theorem 5.1]{ref1} now implies that $A$ has the nSSP.

\end{proof}

\section{Necessity: a deepest spectral collision}


\begin{proposition}

If (H) fails, then $A$ does not have the nSSP.

\end{proposition}

\begin{proof}

Choose a deepest vertex $v$ at which (H) fails. Two children $c,d$ have a common root $\lambda\in\mathbb C$ of $F_c$ and $F_d$. All vertices strictly below $v$ satisfy (H), so Lemma 4.1 applies to both child subtrees and their transposes.

Choose right $\lambda$-eigenvectors $z_c,z_d$ of $B_c,B_d$ and left eigenvectors $y_c,y_d$, written as right eigenvectors of $B_c^{\mathsf T},B_d^{\mathsf T}$. Normalize all four child-root coordinates to $1$. Let

\[
p_c=A_{vc},\quad p_d=A_{vd},
\qquad
q_c=A_{cv},\quad q_d=A_{dv}.
\]

Define whole-tree vectors, zero at $v$, its ancestors, and every unused branch, by

\[
u_+|_{T_c}=p_dz_c,
\qquad
u_+|_{T_d}=-p_cz_d,
\tag{7.1}
\]

and

\[
w_+|_{T_c}=q_dy_c,
\qquad
w_+|_{T_d}=-q_cy_d.
\tag{7.2}
\]

The coefficients cancel the row and column equations at $v$, respectively. Hence

\[
Au_+=\lambda u_+,
\qquad
A^{\mathsf T}w_+=\lambda w_+.
\tag{7.3}
\]

Each zero-diagonal child subtree is bipartite. Let $S$ be the diagonal sign matrix on the two used subtrees, with sign $+1$ at both child roots and alternating by depth. Put

\[
u_-=Su_+,
\qquad
w_-=Sw_+.
\]

Then $u_-,w_-$ are right and left eigenvectors for $-\lambda$. Consequently,

\[
Y=u_+w_+^{\mathsf T}+u_-w_-^{\mathsf T}
\tag{7.4}
\]

commutes with $A$.

For adjacent vertices $i,j$ within either used child subtree, their bipartition signs satisfy $s_is_j=-1$, so

\[
Y_{ij}=(u_+)_i(w_+)_j(1+s_is_j)=0.
\]

Every edge incident with $v$ has one zero vector coordinate, and every other prescribed edge lies outside the support or inside a used subtree. The root-loop position is also zero. Hence $Y$ vanishes at every prescribed nonzero position of $A$.

At the child root $c$,

\[
Y_{cc}=2p_dq_d=2\beta_{vd}\in\mathbb R\setminus\{0\}.
\tag{7.5}
\]

Thus $Y$ is nonzero and $\operatorname{Re}Y$ is a nonzero real commuting matrix. Let

\[
X=(\operatorname{Re}Y)^{\mathsf T}.
\]

The support of $A$ is bidirected, so transposition preserves the pattern-zero condition. Therefore $X\ne0$ satisfies (2.1), and $A$ fails the nSSP.

The argument includes $\lambda=0$: the sign transform may return the same eigenvectors, but the edge cancellation and (7.5) remain valid. It also includes a repeated or nonreal common root because only left and right eigenvectors, not diagonalizability, are used.

\end{proof}

\section{The classification theorem and consequences}


\begin{theorem}

For the root-loop bidirected tree matrix $A$ in Section 1, the following are equivalent:

\begin{enumerate}
\item $A$ has the nSSP;
\item for every vertex $v$, the child-subtree polynomials $\{F_c:c\in C(v)\}$ are pairwise coprime.
\end{enumerate}


\end{theorem}

\begin{proof}

Proposition 6.2 proves sufficiency. Proposition 7.1 proves the contrapositive of necessity.

\end{proof}

\begin{corollary}[NEB obstruction]

If $A$ has the nSSP, then $T$ is nearly evenly branching at $r$ in the recursive sense of \cite{ref4}.

\end{corollary}

\begin{proof}

Every odd-order zero-diagonal child subtree has $F_c(0)=0$. Pairwise coprimality therefore permits at most one odd-order child at each vertex. If $|T_v|$ is even, the sum of child orders is odd, so exactly one child has odd order. If $|T_v|$ is odd, the sum is even, so no child has odd order. Recursing gives the NEB condition.

The converse is false in general: NEB removes the forced common factor $x$, but two sibling subtrees may still share a nonzero or nonreal eigenvalue.

\end{proof}

\begin{corollary}[recursive certificate]

When the matrix entries lie in an effective exact field---for example $\mathbb Q$, an explicitly represented algebraic-number field, or a symbolic field with exact zero testing---the nSSP status can be decided bottom-up by exact polynomial arithmetic: compute every $F_v$ from (3.1), then compute the pairwise gcds at each sibling set. A nonconstant gcd supplies a spectral-collision certificate and Proposition 7.1 supplies a centralizer witness over a splitting field; if every gcd is $1$, Lemma 4.1 and Propositions 6.1--6.2 certify the nSSP. This is not asserted as a numerical algorithm for arbitrary inexact real inputs.

\end{corollary}

\begin{remark}

A repeated root inside one $F_c$ does not obstruct the nSSP unless the same factor occurs in a sibling polynomial. Thus the theorem is not a simple-spectrum criterion.

\end{remark}

\begin{remark}

For a path rooted at an endpoint, every sibling set has at most one member and the condition is vacuous. For a root-loop double path it becomes separation of the two arms \cite{ref12}. For a root-loop spider it becomes pairwise coprimality of the arm polynomials \cite{ref7}.

\end{remark}

\begin{remark}

The zero nonroot diagonal hypothesis cannot simply be removed from the sufficient direction. On the endpoint-rooted path $0-1-2$, let

\[
A=
\begin{pmatrix}
1&1&0\\
1&0&1\\
0&1&1
\end{pmatrix},
\qquad
X=A^2-A-I=
\begin{pmatrix}
0&0&1\\
0&1&0\\
1&0&0
\end{pmatrix}.
\]

Every sibling set has at most one member, so the recursive coprimality condition is vacuous. Nevertheless, $X=X^{\mathsf T}\ne0$, $A\circ X=0$, and $AX^{\mathsf T}=X^{\mathsf T}A$ because $X$ is a polynomial in $A$. Thus $A$ does not have the nSSP. This refutes the same sufficient criterion for arbitrary fixed nonroot diagonal entries; it does not assert that every additional nonroot loop destroys the nSSP.

\end{remark}

\section{Exact computational calibration}


The proof is independent of computation. Three dependency-free Python programs provide definition-level falsification controls:

\begin{itemize}
\item \href{verification/verify_recursive_tree_criterion.py}{recursive criterion calibration};
\item \href{verification/verify_collision_witness.py}{explicit collision witnesses};
\item the \href{verification/verify_minimal_counterexample.py}{exact three-vertex boundary certificate} from Remark 8.6.
\end{itemize}


The standard calibration command is

\begin{Verbatim}[breaklines=true,breakanywhere=true,fontsize=\small]
python -X utf8 verification/verify_recursive_tree_criterion.py --max-vertices 7 --samples 2
\end{Verbatim}


It covers all $1,2,6,24,120,720$ increasing-rooted parent tuples on $n=2,\ldots,7$ vertices, with two signed directed-weight samples each, for 1746 cases total.

\begin{itemize}
\item all 366 recursively coprime cases have full verification rank modulo $1000003$;
\item none of the 1380 recursively noncoprime cases has modular full rank;
\item six targeted cases independently reproduce the subtree polynomials by a Newton--trace implementation and use exact rational elimination where a rank loss is asserted.
\end{itemize}


The targeted controls include nested branch collisions, two distinct branch vertices, negative edge products, a single subtree with repeated complex roots, and a changed directional splitting with fixed products.

The explicit-witness command is

\begin{Verbatim}[breaklines=true,breakanywhere=true,fontsize=\small]
python -X utf8 verification/verify_collision_witness.py
\end{Verbatim}


It constructs rational pattern-zero centralizers directly for a deepest shared zero root and a deepest shared factor $x^2-1$. Both commute exactly with $A$, vanish on every prescribed position, and are nonzero.

The boundary-certificate command is

\begin{Verbatim}[breaklines=true,breakanywhere=true,fontsize=\small]
python -X utf8 verification/verify_minimal_counterexample.py
\end{Verbatim}

It reconstructs $X=A^2-A-I$, checks $X=X^{\mathsf T}$ and the transpose commutation equation directly, recomputes the verification rank as $2/3$ over $\mathbb Q$, and checks two destructive controls of rank $3/3$.

Modular full rank proves full rank over $\mathbb Q$ and $\mathbb R$ for that integer instance. Modular rank loss alone does not prove rational rank loss. Random and finite cases serve only to catch false conjectures and implementation errors; they do not prove Theorem 8.1.

\section{Relation to prior work}


Duarte's inverse spectral construction for acyclic Hermitian matrices \cite{ref8} and the Monfared--Shader Jacobian formulation \cite{ref3} recursively impose strict interlacing between a tree matrix and a vertex-deleted principal submatrix. Their Duarte property yields a strong-Arnold-like centralizer lemma and a nonsingular joint spectral-data Jacobian. For real symmetric tree matrices, recursive strict interlacing is closely equivalent to disjoint sibling spectra together with cyclic child roots. The positive-product subfamily of the present theorem is real diagonally similar to that symmetric setting and is not claimed as new.

Monfared--Mallik \cite{ref4} develop the skew-symmetric Duarte analogue, characterize the NEB tree obstruction, and prove a corresponding Jacobian result for distinct purely imaginary interlacing data. Corollary 8.2 is therefore a compatibility boundary, not a new NEB mechanism. The present theorem differs by deciding every fixed matrix with mixed products and by allowing nonreal and defective subtree spectra.

Generalized-star and linear-tree inverse spectral results \cite{ref13,ref14}
also use branch spectra and multiplicity constraints.  Their quantifiers concern
existence or realization of matrices with prescribed spectral data, rather
than the nSSP of every fixed nonsymmetric root-loop matrix; they do not state
the recursive iff proved here.

Monfared \cite{ref5} proves an existence theorem for not-necessarily-symmetric matrices with a prescribed graph and distinct real or conjugate-pair spectrum, under the hypothesis that the graph's matching number is at least the number of prescribed nonreal conjugate pairs. That work does not classify the nSSP of a fixed rooted-tree matrix.

Fallat, Hall, Lin, and Shader \cite{ref2} formulate the nSSP and its bifurcation/superpattern framework. Arav et al. \cite{ref11} characterize the same similarity transversality through the STP and develop its pattern-level consequences. Catral, Fallat, Gupta, and Lin \cite{ref1} later prove the exact characteristic-Jacobian equivalence used in Proposition 6.2. These are essential prior bridges, but none states a fixed rooted-tree recursive-coprimality criterion.

Koljan\v{c}i\'{c}--Oblak \cite{ref6} study which directed graph patterns allow or require nSSP and classify many loop assignments on double paths. Their path results contain endpoint-rooted path special cases at the pattern level, but they do not state the arbitrary fixed-tree hierarchy in Theorem 8.1. The author's earlier public Beta \cite{ref12} proves a spectral-separation criterion for symmetric single-loop double paths and uses it in a pattern-level classification. Its fixed-matrix criterion is recovered by Theorem 8.1, while its classification of arbitrary loop assignments is not reproved here.

The author's two preceding results form the progression double path \cite{ref12}, root-loop spider \cite{ref7}, and arbitrary rooted tree. Theorem 8.1 shows that the spider pairwise-coprimality mechanism is the first level of a fully recursive tree classification. The two earlier fixed-matrix criteria remain useful precursors with simpler proofs and verifiers, but the general theorem strictly subsumes those fixed-matrix statements; the separate pattern-level loop classification in \cite{ref12} remains independent.

The public sources above were checked through 2026-08-19. No direct arbitrary mixed-product fixed-tree nSSP iff was found. Search non-detection does not establish global priority, and any public version should retain that limitation.

\section{Limits and further questions}


The proof uses each frozen assumption.

\begin{itemize}
\item The unique nonzero root loop separates the parity components of $\chi_A=F_r-aQ_r$.
\item Zero nonroot diagonal entries provide the $\lambda/-\lambda$ symmetry used by the failure witness; Remark 8.6 shows that deleting this hypothesis while retaining the same sufficient criterion is already false on three vertices.
\item Nonzero weights make the product differential onto and prevent a branch from decoupling.
\item Bidirected support is needed for two-way products and for the transpose of the witness to preserve pattern zeros.
\item A tree makes the bottom-up factorization and deepest-collision localization exact.
\end{itemize}


Several extensions are therefore separate problems rather than routine corollaries:

\begin{enumerate}
\item several nonzero loops on one tree;
\item missing reverse arcs or zero edge weights;
\item unicyclic or block-tree support;
\item a pattern-level characterization of which rooted-tree sign patterns allow the criterion to be realized;
\item efficient construction of a rational witness when the common factor is irreducible over $\mathbb Q$.
\end{enumerate}


The current paper should not absorb these questions before release. Their value can be assessed in separate bounded loops.

\section{Reproducibility and evidence boundary}


The public research package contains:

\begin{itemize}
\item this manuscript PDF and its corresponding \href{nssp-root-loop-trees.tex}{LaTeX source};
\item an \href{README.md}{entry README} stating the theorem, scope, status, and reproduction commands;
\item the three dependency-free verification programs linked in Section 9 and their \href{verification/README.md}{verification guide};
\item the \href{SHA256SUMS.txt}{SHA-256 manifest} and the entry-specific license notices.
\end{itemize}

The finite tests serve only to catch implementation errors, exercise boundary cases, and prevent hidden symmetric or simple-spectrum assumptions. The theorem is supported by the proof in Sections 3--8. No solver timeout, modular rank loss, random sample, or finite enumeration is used as a general theorem.

\section{AI-assisted research and writing disclosure}


This work was developed with AI assistance for literature search, symbolic derivation, proof red-teaming, code generation, exact computational checks, and manuscript drafting. Carptopus is the responsible author and is accountable for the mathematical claims, source selection, software, and public versions. OpenAI Codex is a research and writing tool, not an author.

\begin{thebibliography}{12}

\bibitem{ref1} M. Catral, S. M. Fallat, H. Gupta, and J. C.-H. Lin,
\emph{The Strong Spectral Property and the Jacobian Method for Weighted Laplacian Matrices},
arXiv:2602.18999v1 (2026), especially Theorem 5.1.
\url{https://arxiv.org/abs/2602.18999}

\bibitem{ref2} S. M. Fallat, H. T. Hall, J. C.-H. Lin, and B. L. Shader,
\emph{The bifurcation lemma for strong properties in the inverse eigenvalue problem of a graph},
Linear Algebra and its Applications 648 (2022), 70--87.
\url{https://arxiv.org/abs/2112.11712}

\bibitem{ref3} K. H. Monfared and B. L. Shader,
\emph{Construction of matrices with a given graph and prescribed interlaced spectral data},
Linear Algebra and its Applications 438 (2013), 4348--4358.
\url{https://doi.org/10.1016/j.laa.2013.01.036}

\bibitem{ref4} K. H. Monfared and S. Mallik,
\emph{Construction of real skew-symmetric matrices from interlaced spectral data, and graph},
Linear Algebra and its Applications 471 (2015), 241--263.
\url{https://arxiv.org/abs/1412.6085}

\bibitem{ref5} K. H. Monfared,
\emph{Existence of a not necessarily symmetric matrix with given distinct eigenvalues and graph},
Linear Algebra and its Applications 527 (2017), 1--11.
\url{https://arxiv.org/abs/1604.02195}

\bibitem{ref6} S. Koljan\v{c}i\'{c} and P. Oblak,
\emph{Combinatorial aspects of the non-symmetric strong spectral property for graphs},
arXiv:2604.21552v1 (2026).
\url{https://arxiv.org/abs/2604.21552}

\bibitem{ref7} Carptopus,
\emph{An exact non-symmetric strong spectral criterion for root-loop spider matrices},
MathExplore Notes, v0.1-beta (19 August 2026), commit \texttt{26d6458}.
\href{https://github.com/Carptopus/MathExplore-Notes/tree/26d6458029d09c59fd52a56ede4840ed7faf7bae/research/nssp-root-loop-spiders}{immutable GitHub source}

\bibitem{ref8} A. L. Duarte,
\emph{Construction of acyclic matrices from spectral data},
Linear Algebra and its Applications 113 (1989), 173--182.
\url{https://doi.org/10.1016/0024-3795(89)90295-4}

\bibitem{ref9} R. A. Horn and C. R. Johnson,
\emph{Matrix Analysis}, second edition, Cambridge University Press, 2013.

\bibitem{ref10} T. Kailath,
\emph{Linear Systems}, Prentice--Hall, 1980.

\bibitem{ref11} M. Arav, F. J. Hall, H. van der Holst, Z. Li, A. Mathivanan,
J. Pan, H. Xu, and Z. Yang,
\emph{Advances on similarity via transversal intersection of manifolds},
Linear Algebra and its Applications 721 (2025), 102--121.
\url{https://doi.org/10.1016/j.laa.2024.05.013}

\bibitem{ref12} Carptopus,
\emph{Spectral separation and the non-symmetric strong spectral property for looped double paths},
MathExplore Notes, v0.1-beta (19 August 2026), commit \texttt{e57ee00}.
\href{https://github.com/Carptopus/MathExplore-Notes/tree/e57ee007380cbc351eb371d8e991303cab5f1e48/research/nssp-looped-double-paths}{immutable GitHub source}

\bibitem{ref13} C. R. Johnson, A. Leal Duarte, and C. M. Saiago,
\emph{Inverse eigenvalue problems and lists of multiplicities of eigenvalues for matrices whose graph is a tree: the case of generalized stars and double generalized stars},
Linear Algebra and its Applications 373 (2003), 311--330.
\url{https://doi.org/10.1016/S0024-3795(03)00582-2}

\bibitem{ref14} C. R. Johnson and T. Wakhare,
\emph{The inverse eigenvalue problem for linear trees},
Discrete Mathematics 345 (2022), Article 112737.
\url{https://doi.org/10.1016/j.disc.2021.112737}

\end{thebibliography}

\end{document}
